\section{Defenses}
\label{sec:defense}

Standard boosting chooses split features by training-loss reduction alone, so features that separate the classes widely in the synthetic corpora, such as duration and the rates, sit near the root although an attacker can shift them by delaying or padding packets. We evaluate malleability-penalized splitting (MPS), which discourages splits on malleable features, and Lie-guided adversarial training (LGAT), which trains on flows generated along the Lie generators.

\subsection{Malleability-penalized tree splitting}

Following the robust split formulations of Chen et al.~\cite{chen2019robust} and Calzavara et al.~\cite{calzavara2020treant}, the split score for feature $j$ is regularized by its physical malleability index.

\begin{definition}[Feature Physical Malleability Index]
The physical malleability index $M_{\text{mal}}(j)$ of feature $j$ is defined as the empirical expectation over the training distribution of the maximum dimensionless variation induced by the Lie Cone $\mathcal{K}_{\text{net}}$:
\begin{equation}
    M_{\text{mal}}(j) = \mathbb{E}_{x \sim \mathcal{D}_{\text{train}}} \left[ \max_{k \in \{1, \dots, 4\}} \frac{\delta s_k |\vec{v}_k(x)_j|}{\sigma_j} \right]
\end{equation}
where $\sigma_j$ is the empirical standard deviation of feature $j$ across the training corpus $\mathcal{D}_{\text{train}}$, and $\delta s_k$ denotes the characteristic single-step perturbation scale for generator $k$ ($\delta s_1 = 10^4\,\mu\text{s}, \delta s_2 = 100\,\text{B}, \delta s_3 = 1\,\text{pkt}, \delta s_4 = 10^3\,\text{B}$).
\end{definition}

The robust split selection objective is
\begin{equation}
    \text{Gain}_{\text{robust}}(j, \theta) = \Delta \mathcal{L}(j, \theta) - \alpha_{\text{mal}} \cdot \overline{\Delta \mathcal{L}} \cdot M_{\text{mal}}(j)
\end{equation}
where $\Delta \mathcal{L}(j, \theta)$ is the impurity gain of split $(j, \theta)$, $\overline{\Delta \mathcal{L}}$ is the baseline mean impurity gain across candidate splits at the node, ensuring dimensional consistency, and $\alpha_{\text{mal}} \ge 0$ is a dimensionless hyperparameter governing the trade-off between clean training loss and malleability robustness.

The experiments do not optimize this objective. The MPS model instead masks duration ($x_6$) and the two rate features ($x_8$, $x_9$) to zero in the training, test and attack inputs, which is the $\alpha_{\text{mal}} \to \infty$ limit for those features. The mask is fixed by hand rather than computed from $M_{\text{mal}}$, and other features are unpenalized. Because the masked columns are constant, trees cannot split on them, so the attacker affects MPS models only through unmasked features.

\subsection{Lie-guided adversarial training}

Standard adversarial training assumes an $L_p$ ball around each input~\cite{madry2018towards}, which in network security generates physically invalid samples that degrade clean accuracy~\cite{hashemi2020enhancing, apruzzese2022modeling}. LGAT instead restricts the inner maximization to the physical flow manifold, solving
\begin{equation}
    \min_{\Theta} \frac{1}{N} \sum_{i=1}^N \max_{u \in \mathcal{U}_{\text{budget}}} \mathcal{L}\left( F_\Theta\left( \mathcal{A}_{\text{cone}}\left( x_i, u \right) \right), y_i \right)
\end{equation}
where $\mathcal{A}_{\text{cone}}(x, u)$ is the flow reached from $x$ along the integral curves of the Lie generators under the control $u$, and $\mathcal{U}_{\text{budget}}$ is the set of nonnegative controls whose cost $\sum_{k=1}^4 c_k \int_0^{T_{\max}} u_k(\tau)\, d\tau$ does not exceed $\Gamma$.

Supplementary Algorithm~S4 gives the general scheme, in which malicious flows are replaced by FFE outputs at step budget $M_{\text{train}} = 5$ in every boosting round, with $\beta = 0.5$. The experiments use a single-round instance: adversarial examples are generated once, against the undefended model, for a random 50\% of the malicious training flows and appended to the training set, and the ensemble uses standard split scoring.

\subsection{Evaluation}

We compared an undefended XGBoost model with three hardened variants on the CSE-CIC-IDS2018 calibrated corpus, each with 60 trees of depth 6 and learning rate 0.1. All models are attacked on the same cohort of 100 flows that the undefended model flags, using FFE at step budget $M=5$ (Table~\ref{tab:defense_evaluation}).

\begin{table}[t]
\caption{Defensive hardening against Flow-Field Evasion on CSE-CIC-IDS2018 (XGBoost; calibrated corpus; cohort $N=100$; attack budget $M=5$; all values in \%). Robust recall is $100 - \text{ASR}$.}\label{tab:defense_evaluation}
\linespread{1}\selectfont
\centering
\begin{tabular}{@{}lcccc@{}}
\toprule
Model & Clean acc. & Clean AUC & Robust recall & FFE ASR \\
\midrule
Undefended    & 99.72 & 99.99 & 5.0  & 95.0 \\
MPS-hardened  & 97.80 & 99.81 & 23.0 & 77.0 \\
LGAT-hardened & 99.25 & 99.97 & 12.0 & 88.0 \\
MPS + LGAT    & 97.30 & 99.69 & 44.0 & 56.0 \\
\bottomrule
\end{tabular}
\end{table}

Combining MPS with LGAT reduces FFE success from 95.0\% to 56.0\% and raises robust recall from 5.0\% to 44.0\%. The clean accuracy cost is 2.42 percentage points (99.72\% to 97.30\%), with clean ROC-AUC of 99.69\%. The attack budget explains the undefended figure: at $M=5$ FFE succeeds on 95.0\% of the cohort, whereas Section~\ref{sec:results} reports 100.0\% at the larger budget. The residual 56.0\% success shows that the defense is partial.
